\section{The full original character and specialization-compatible projectors}
\label{sec:gct-full-character}

The original row shape in this section is $\nu=(7,5,3)$, its column
lengths are $(3,3,3,2,2,1,1)$, and its column-count shape is
$[0,3,2,2]$. None of these data is replaced by a smaller partition.
Set $A=\mathbb Z[q,q^{-1}]$, $K=\mathbb Q(q)$, and
$X_A=A\,\mathrm{pNSTC}(\nu)$ inside the source quotient
\[
 X=\check X_\nu=\check Y_{(3,3,3,2,2,1,1)}/
 \sum_{i=1}^{6}\check Y_{\triangleright^i(3,3,3,2,2,1,1)}.
\]
Write $[t]=(q^t-q^{-t})/(q-q^{-1})$ for every integer $t$,
with $[0]=0$ and $[-t]=-[t]$. The source parameter is retained.

We use the actual two-row theorem of the April 2013 GCT-IV source,
whose TeX labels are \texttt{t main canonical basis},
\texttt{p highest weight straightened NST}, and
\texttt{c basis correct size}. Its hypothesis is
$\dim V=\dim W=2$, with arbitrary $\nu$ of length at most four;
it is not a theorem constructing all conjectured nonstandard tensor
fat cells. At this scope the positive honest classes are an $A$-basis,
the divided Chevalley generators preserve that lattice, and the
crystal cells have the highest weights obtained by the source's
straightened-class rule. These are facts about the quotient just
displayed. Crystal-highest class representatives are not being
identified with simultaneous Chevalley kernels in $X$.

\subsection{An exhaustive count in the original column shape}

An original straightened highest class is specified by its invariant
record $(0,i_3,i_2,i_1)$ and invariant-free columns
\[
 (123)^A(124)^B(134)^C(12)^D(13)^E(23)^G(1)^f.
\]
Here powers count consecutive columns, not commutative products.
All signs and scalings remain those of the positive source class:
an underlying tabloid of degree $h=i_3+i_2+i_1$ has the source
factor $(-1/[2])^h$, followed by its prescribed orientation sign.
Counting a class does not change this scalar.
The inserted height-three invariant is $(234)(123)$,
the prescribed height-two invariant is $(24)(13)$, and the
height-one invariant is $(4)(1)$. They are positioned by the
source's straightening convention; their respective two-sided
weights are $(3,3;3,3)$, $(2,2;2,2)$, and $(1,1;1,1)$.
Thus the complete list of constraints in the retained shape is
\begin{gather}
 i_3,i_2,i_1\in\{0,1\},\quad
 A+B+C=3-2i_3,\quad D+E+G=2-2i_2,\quad f=2-2i_1,
 \label{eq:gct-honest-shape}\\
 A,B,C,D,E,G,f\geq0,\quad BC=0,\quad
 D>0\Longrightarrow C=0,\quad E>0\Longrightarrow B\leq1,
 \nonumber\\
 G\leq1,\quad G=1\Longrightarrow f>0,\quad
 B\leq f+E,\quad C\leq f+D.\label{eq:gct-honest-inequalities}
\end{gather}
The source's highest-class classification and invariant insertion
give one positive honest highest class for each solution, and its
unique straightened representative recovers these integers. Hence
there is neither an omitted class nor multiplicity from choosing
different representatives of one class.

Put $k=3i_3+2i_2+i_1$. The original alphabet
$1=(1,1)$, $2=(1,2)$, $3=(2,1)$, $4=(2,2)$ gives, by counting
each letter in these columns,
\begin{align}
 a&=2A+2B+C+2D+E+G+f+k,\nonumber\\
 b&=2A+B+2C+D+2E+G+f+k,\label{eq:gct-highest-letter-weights}\\
 \lambda&=(a,15-a),\qquad \mu=(b,15-b).\nonumber
\end{align}
The invariants add the displayed $k$ to both coordinates of both
weights. Consequently the total degree stays fifteen on each side.
The inequalities ensure $a,b\geq8$.

Here is the entire finite enumeration. Each row is one solution of
\eqref{eq:gct-honest-shape}--\eqref{eq:gct-honest-inequalities};
the last column is $(2a-14)(2b-14)$.
\begin{center}\small
\begin{tabular}{ccccccr}
\toprule
$(i_3,i_2,i_1)$ & $(A,B,C)$ & $(D,E,G)$ & $f$ & $\lambda$ & $\mu$ & dimension\\
\midrule
$(0,0,0)$ & $(1,0,2)$ & $(0,1,1)$ & $2$ & $(8,7)$ & $(11,4)$ & 16\\
$(0,0,0)$ & $(1,0,2)$ & $(0,2,0)$ & $2$ & $(8,7)$ & $(12,3)$ & 20\\
$(0,0,0)$ & $(1,2,0)$ & $(1,0,1)$ & $2$ & $(11,4)$ & $(8,7)$ & 16\\
$(0,0,0)$ & $(1,2,0)$ & $(2,0,0)$ & $2$ & $(12,3)$ & $(8,7)$ & 20\\
$(0,0,0)$ & $(2,0,1)$ & $(0,1,1)$ & $2$ & $(9,6)$ & $(11,4)$ & 32\\
$(0,0,0)$ & $(2,0,1)$ & $(0,2,0)$ & $2$ & $(9,6)$ & $(12,3)$ & 40\\
$(0,0,0)$ & $(2,1,0)$ & $(0,1,1)$ & $2$ & $(10,5)$ & $(10,5)$ & 36\\
$(0,0,0)$ & $(2,1,0)$ & $(0,2,0)$ & $2$ & $(10,5)$ & $(11,4)$ & 48\\
$(0,0,0)$ & $(2,1,0)$ & $(1,0,1)$ & $2$ & $(11,4)$ & $(9,6)$ & 32\\
$(0,0,0)$ & $(2,1,0)$ & $(1,1,0)$ & $2$ & $(11,4)$ & $(10,5)$ & 48\\
$(0,0,0)$ & $(2,1,0)$ & $(2,0,0)$ & $2$ & $(12,3)$ & $(9,6)$ & 40\\
$(0,0,0)$ & $(3,0,0)$ & $(0,1,1)$ & $2$ & $(10,5)$ & $(11,4)$ & 48\\
$(0,0,0)$ & $(3,0,0)$ & $(0,2,0)$ & $2$ & $(10,5)$ & $(12,3)$ & 60\\
$(0,0,0)$ & $(3,0,0)$ & $(1,0,1)$ & $2$ & $(11,4)$ & $(10,5)$ & 48\\
$(0,0,0)$ & $(3,0,0)$ & $(1,1,0)$ & $2$ & $(11,4)$ & $(11,4)$ & 64\\
$(0,0,0)$ & $(3,0,0)$ & $(2,0,0)$ & $2$ & $(12,3)$ & $(10,5)$ & 60\\
$(0,0,1)$ & $(2,1,0)$ & $(0,2,0)$ & $0$ & $(9,6)$ & $(10,5)$ & 24\\
$(0,0,1)$ & $(2,1,0)$ & $(1,1,0)$ & $0$ & $(10,5)$ & $(9,6)$ & 24\\
$(0,0,1)$ & $(3,0,0)$ & $(0,2,0)$ & $0$ & $(9,6)$ & $(11,4)$ & 32\\
$(0,0,1)$ & $(3,0,0)$ & $(1,1,0)$ & $0$ & $(10,5)$ & $(10,5)$ & 36\\
$(0,0,1)$ & $(3,0,0)$ & $(2,0,0)$ & $0$ & $(11,4)$ & $(9,6)$ & 32\\
$(0,1,0)$ & $(1,0,2)$ & $(0,0,0)$ & $2$ & $(8,7)$ & $(10,5)$ & 12\\
$(0,1,0)$ & $(1,2,0)$ & $(0,0,0)$ & $2$ & $(10,5)$ & $(8,7)$ & 12\\
$(0,1,0)$ & $(2,0,1)$ & $(0,0,0)$ & $2$ & $(9,6)$ & $(10,5)$ & 24\\
\bottomrule
\end{tabular}
\end{center}

\begin{center}\small
\begin{tabular}{ccccccr}
\toprule
$(i_3,i_2,i_1)$ & $(A,B,C)$ & $(D,E,G)$ & $f$ & $\lambda$ & $\mu$ & dimension\\
\midrule
$(0,1,0)$ & $(2,1,0)$ & $(0,0,0)$ & $2$ & $(10,5)$ & $(9,6)$ & 24\\
$(0,1,0)$ & $(3,0,0)$ & $(0,0,0)$ & $2$ & $(10,5)$ & $(10,5)$ & 36\\
$(0,1,1)$ & $(3,0,0)$ & $(0,0,0)$ & $0$ & $(9,6)$ & $(9,6)$ & 16\\
$(1,0,0)$ & $(0,0,1)$ & $(0,1,1)$ & $2$ & $(8,7)$ & $(10,5)$ & 12\\
$(1,0,0)$ & $(0,0,1)$ & $(0,2,0)$ & $2$ & $(8,7)$ & $(11,4)$ & 16\\
$(1,0,0)$ & $(0,1,0)$ & $(0,1,1)$ & $2$ & $(9,6)$ & $(9,6)$ & 16\\
$(1,0,0)$ & $(0,1,0)$ & $(0,2,0)$ & $2$ & $(9,6)$ & $(10,5)$ & 24\\
$(1,0,0)$ & $(0,1,0)$ & $(1,0,1)$ & $2$ & $(10,5)$ & $(8,7)$ & 12\\
$(1,0,0)$ & $(0,1,0)$ & $(1,1,0)$ & $2$ & $(10,5)$ & $(9,6)$ & 24\\
$(1,0,0)$ & $(0,1,0)$ & $(2,0,0)$ & $2$ & $(11,4)$ & $(8,7)$ & 16\\
$(1,0,0)$ & $(1,0,0)$ & $(0,1,1)$ & $2$ & $(9,6)$ & $(10,5)$ & 24\\
$(1,0,0)$ & $(1,0,0)$ & $(0,2,0)$ & $2$ & $(9,6)$ & $(11,4)$ & 32\\
$(1,0,0)$ & $(1,0,0)$ & $(1,0,1)$ & $2$ & $(10,5)$ & $(9,6)$ & 24\\
$(1,0,0)$ & $(1,0,0)$ & $(1,1,0)$ & $2$ & $(10,5)$ & $(10,5)$ & 36\\
$(1,0,0)$ & $(1,0,0)$ & $(2,0,0)$ & $2$ & $(11,4)$ & $(9,6)$ & 32\\
$(1,0,1)$ & $(0,1,0)$ & $(0,2,0)$ & $0$ & $(8,7)$ & $(9,6)$ & 8\\
$(1,0,1)$ & $(0,1,0)$ & $(1,1,0)$ & $0$ & $(9,6)$ & $(8,7)$ & 8\\
$(1,0,1)$ & $(1,0,0)$ & $(0,2,0)$ & $0$ & $(8,7)$ & $(10,5)$ & 12\\
$(1,0,1)$ & $(1,0,0)$ & $(1,1,0)$ & $0$ & $(9,6)$ & $(9,6)$ & 16\\
$(1,0,1)$ & $(1,0,0)$ & $(2,0,0)$ & $0$ & $(10,5)$ & $(8,7)$ & 12\\
$(1,1,0)$ & $(0,0,1)$ & $(0,0,0)$ & $2$ & $(8,7)$ & $(9,6)$ & 8\\
$(1,1,0)$ & $(0,1,0)$ & $(0,0,0)$ & $2$ & $(9,6)$ & $(8,7)$ & 8\\
$(1,1,0)$ & $(1,0,0)$ & $(0,0,0)$ & $2$ & $(9,6)$ & $(9,6)$ & 16\\
$(1,1,1)$ & $(1,0,0)$ & $(0,0,0)$ & $0$ & $(8,7)$ & $(8,7)$ & 4\\
\bottomrule
\end{tabular}
\end{center}


For a fully specified exhaustive procedure, choose the eight invariant
records, then $A=0,\ldots,3-2i_3$,
$B=0,\ldots,3-2i_3-A$, and set $C=3-2i_3-A-B$.
Choose $D=0,\ldots,2-2i_2$ and
$E=0,\ldots,2-2i_2-D$, and set $G=2-2i_2-D-E$.
Reject exactly the failed inequalities above, set $f=2-2i_1$,
and evaluate \eqref{eq:gct-highest-letter-weights}.
Every nonnegative solution is visited exactly once because all its
free coordinates are visited exactly once. This proves the enumeration
without a numerical interpolation or a guess about saturation.

Let $g_{ab}$ be the number of these solutions with highest first
coordinates $(a,b)$. Rows and columns in the following matrix have
the order $12,11,10,9,8$:
\begin{equation}
 (g_{ab})=
 \begin{pmatrix}
 0&0&1&1&1\\0&1&2&3&2\\1&2&4&4&3\\
 1&3&4&4&2\\1&2&3&2&1
 \end{pmatrix}.
 \label{eq:gct-full-multiplicity-matrix}
\end{equation}
All other highest weights have multiplicity zero. The sum of the
entries is $48$, and
\begin{equation}
 \dim_K X=\sum_{a,b=8}^{12}g_{ab}(2a-14)(2b-14)=1260.
 \label{eq:gct-full-dimension}
\end{equation}

\subsection{The full character and the classical specialization map}

Let $x_1,x_2$ and $y_1,y_2$ record the two original torus weights.
The complete character, retaining every coordinate, is
\begin{equation}
 \operatorname{ch}X=
 \sum_{a,b=8}^{12}g_{ab}
 \left(\sum_{i=0}^{2a-15}x_1^{a-i}x_2^{15-a+i}\right)
 \left(\sum_{j=0}^{2b-15}y_1^{b-j}y_2^{15-b+j}\right).
 \label{eq:gct-full-character-polynomial}
\end{equation}
Each crystal component is a rectangle with one basis vector at
every listed weight, so the source highest-class count proves this
formula for the original quotient, not for an unrelated space of
the same dimension.

For reference, the dimensions $w_{uv}$ of the dominant weight
spaces $(u,15-u;v,15-v)$, with rows and columns again in order
$12,11,10,9,8$, are
\begin{equation}
 (w_{uv})=
 \begin{pmatrix}
 0&0&1&2&3\\0&1&4&8&11\\1&4&11&19&25\\
 2&8&19&31&39\\3&11&25&39&48
 \end{pmatrix}.
 \label{eq:gct-full-dominant-weight-matrix}
\end{equation}
The other weights satisfy $w_{uv}=w_{15-u,v}=w_{u,15-v}$,
and vanish outside $3\leq u,v\leq12$. Formula
\eqref{eq:gct-full-character-polynomial} proves these reflection
identities, since each of its two intervals is invariant under the
corresponding reflection. On the dominant chamber the exact inverse
to rectangle summation is
\[
 g_{ab}=w_{ab}-w_{a+1,b}-w_{a,b+1}+w_{a+1,b+1}.
\]
Indeed $w_{ab}=\sum_{c\geq a,d\geq b}g_{cd}$ there, and both
successive finite differences cancel precisely the terms outside
$(c,d)=(a,b)$.

The source specialization is a diagram of actual morphisms
\[
 X\ \longleftarrow\ X_A\ \longrightarrow\
 X_A/(q-1)X_A\ \longrightarrow\
 \mathbb Q\otimes_{\mathbb Z}(X_A/(q-1)X_A)
 \xrightarrow{\ \sigma\ }
 S_{(7,5,3)}(V_{\mathbb Q}\otimes W_{\mathbb Q}).
\]
The first arrow is inclusion followed by scalar extension; the second
is the quotient. We make explicit the source's two integral presentations.
With $Y_A^S$ its rescaled tabloid lattice and $M_A^S$ the sum of its
six specified integral relation lattices, the source first forms
$X'_A=Y_A^S/M_A^S$. The inclusion
$M_A^S\subseteq (K M_A^S)\cap Y_A^S$ gives the actual surjection
\[
 s_X:X'_A\longrightarrow
 Y_A^S/((K M_A^S)\cap Y_A^S)=X_A,\qquad [y]\longmapsto[y].
\]
The source's specialization theorem is stated first for $X'_A$.
Its remark at TeX lines 5901--5908 also asserts an integral
identification, but that stronger assertion is not used here:
an integral relation presentation can retain torsion that the
honest image lattice does not retain. We now prove exactly the
rational specialization map required for the displayed diagram.
Let $\mathbb Q_1$ denote $\mathbb Q$ as an $A$-algebra through
$q\mapsto1$. Right exactness gives a surjective module map
\[
 \overline s_X=\mathbb Q_1\otimes_A s_X:
 \mathbb Q_1\otimes_A X'_A\twoheadrightarrow
 \mathbb Q_1\otimes_A X_A.
\]
The original contextual relation graph has $36864$ tabloid
vertices; its specialization at $q=1$, $p=[2]=2$, has relation
rank $35604$, so its quotient has dimension $1260$ over
$\mathbb Q$. All original rescaling factors are signed powers
of $p$, hence units over $\mathbb Q_1$. They do not change
that rational row span. This is the explicit source relation
specialization computation, recorded in
\texttt{specialization\_bridge.json} with its graph hash and
the original scaled-image check. Equivalently it is the
rational dimension statement in the source's two-row theorem.
Thus the domain of $\overline s_X$ has dimension $1260$.
Its codomain has dimension $1260$ because the original honest
lattice has its proved $1260$-element $A$-basis.

The inverse can be given directly, without an inverse to $s_X$
over $A$. For each original positive honest basis vector $b$,
choose its displayed positive rescaled tabloid representative
$\widetilde b\in Y_A^S$. Define
\[
 j:\mathbb Q_1\otimes_A X_A\longrightarrow
 \mathbb Q_1\otimes_A X'_A,\qquad
 1\otimes b\longmapsto1\otimes[\widetilde b].
\]
This is a well-defined linear map on a basis, and
$\overline s_X j=1$ by the quotient definitions. These $1260$
lifted vectors are independent, since their images are the
honest basis; the dimension calculation proves that they are
also a basis of the domain. It follows that
$j\overline s_X=1$. As the inverse of the surjective module
map $\overline s_X$, $j$ intertwines the specialized original
actions. In particular its value is independent of the choice
of representative. This proves the exact rational isomorphism
and makes no injectivity assertion about $s_X$ over $A$.

The last arrow in the displayed diagram is $\sigma=\tau j$,
where $\tau$ is the source theorem's rational Schur-quotient
map, described intrinsically as follows. The nonstandard exterior
algebra is the quotient of the tensor algebra on the four original
letters by its nonstandard symmetric-square relations. At $q=1$
these relations are the ordinary symmetric-square relations. The
identity on the four letters therefore induces an isomorphism of
the specialized exterior algebra with the ordinary exterior algebra.
Apply its degree maps in the original seven tensor factors. The
source's adjacent-column relation spaces specialize to the adjacent
Schur relation spaces. Thus the tensor map carries the sum of the
six relation spaces onto the sum of the six classical relation
spaces, and it induces $\tau$ on their rational quotients. Applying the
inverse tensor map gives an inverse on these quotients. Naturality
on the four-letter tensor algebra shows that these maps intertwine
the $\mathfrak{gl}_2\oplus\mathfrak{gl}_2$ actions and their
polynomial $\mathrm{GL}_2\times\mathrm{GL}_2$ actions over
$\mathbb Q$. This is the relation proved in the source's
specialization theorem; it does not assert specialization of arbitrary
elements of $K\otimes_A X_A$ that have poles at $q=1$.

In particular
\[
 \operatorname{ch}X=s_{(7,5,3)}
 (x_1y_1,x_1y_2,x_2y_1,x_2y_2).
\]
There is also a separate exact verification of the entire polynomial.
Enumerate all weakly increasing rows of lengths $7,5,3$ on
$\{1,2,3,4\}$, retaining a triple of rows exactly when its columns
strictly increase. This gives all $1260$ semistandard tableaux
exactly once. If their contents are $(c_1,c_2,c_3,c_4)$, their
monomial is
\[
 x_1^{c_1+c_2}x_2^{c_3+c_4}
 y_1^{c_1+c_3}y_2^{c_2+c_4}.
\]
The resulting weight multiplicities equal
\eqref{eq:gct-full-character-polynomial} at every weight.
The separate checker
\texttt{independent\_check/check\_character.py} additionally evaluates the Jacobi--Trudi
determinant in exact integer polynomials. The source's set bijection
between honest classes and ordinary semistandard tableaux is not
being used as a weight-preserving linear map; $\sigma$ is the
module morphism that establishes the classical connection.

\subsection{Explicit intertwiners into the actual source module}

We prove the needed splitting and then give maps whose input
operators are the full original source generators. For
$U_q(\mathfrak{sl}_2)$ use
$KEK^{-1}=q^2E$, $KFK^{-1}=q^{-2}F$ and
$EF-FE=(K-K^{-1})/(q-q^{-1})$.
An integer $h$ denotes a weight on which $K$ acts by $q^h$.

\begin{lemma}\label{lem:gct-explicit-string-splitting}
Every finite-dimensional type-one integrable $K$-module for this
algebra is the direct sum of its highest-weight strings. A string
from a vector $z$ with $Ez=0$ and $Kz=q^n z$ has $n\geq0$,
basis $F^{(i)}z$, $0\leq i\leq n$, and actions
\[
 F F^{(i)}z=[i+1]F^{(i+1)}z,\qquad
 E F^{(i)}z=[n-i+1]F^{(i-1)}z.
\]
\end{lemma}
\begin{proof}
Repeatedly applying the displayed commutator gives
$EF^jz=[j][n-j+1]F^{j-1}z$. This follows by induction:
moving $E$ past one additional $F$ adds the next Cartan term,
and the finite Laurent identity
$[j+1][n-j]-[j][n-j+1]=[n-2j]$ supplies precisely that term.
Let $d$ be the last nonzero power of $F$ on $z$; it exists by
integrability. Applying $E$ to $F^{d+1}z=0$ gives
$[d+1][n-d]F^dz=0$. Over $K$, $[t]=0$ for an integer $t$
only when $t=0$, so $d=n\geq0$. Distinct string vectors have
distinct weights and are independent. Dividing by $[i]!$ yields
the claimed formulas.

For completeness, splitting is proved by induction on dimension.
Choose the largest weight $n$ occurring in a nonzero module.
Its entire weight space consists of highest vectors. The strings
from a basis of that space are independent: a relation at weight
$n-2i$ can be raised by $E^i$, which multiplies each top vector
by the same nonzero scalar. They span a submodule $N$, a sum
of copies of the length-$n$ string. The quotient has smaller
dimension and no weight $n$, so by induction it is a sum of
strings with highest weights $0\leq m<n$.
Lift each quotient highest vector to a vector $z$ of weight $m$.
Then $Ez\in N_{m+2}$. On every copy of the length-$n$ string,
$E:N_m\to N_{m+2}$ is surjective whenever the target occurs:
its coefficient in the formulas above is nonzero. Subtract a
preimage in $N_m$ to obtain a highest lift. Its length-$m$
string maps isomorphically onto the chosen quotient string by
the preceding formulas. All these lifted strings together map
isomorphically to the quotient, so their sum is disjoint from
$N$ and splits the original module. This completes the induction.
\end{proof}

For $a,b\in\{8,9,10,11,12\}$ put $n=2a-15$, $m=2b-15$ and
define the actual source subspace
\[
 H_{ab}=X_{(a,15-a;b,15-b)}\cap\ker E_V\cap\ker E_W.
\]
Let $L_{(a,15-a)}$ be the standard quantum $\mathfrak{gl}_2$
module with basis $e_i$, $0\leq i\leq n$, of weight
$(a-i,15-a+i)$, with $Fe_i=[i+1]e_{i+1}$ and
$Ee_i=[n-i+1]e_{i-1}$; the out-of-range terms are zero.
Define $L_{(b,15-b)}$ in the same way, with basis $f_j$.
In the next tensor product $H_{ab}$ is used as a multiplicity
vector space with trivial quantum-group action, not with an
action inherited from its embedding in $X$. Then the map
\begin{equation}
 \Phi:\bigoplus_{a,b=8}^{12}
 L_{(a,15-a)}\otimes L_{(b,15-b)}\otimes H_{ab}
 \longrightarrow X,\qquad
 e_i\otimes f_j\otimes z\longmapsto F_V^{(i)}F_W^{(j)}z
 \label{eq:gct-actual-string-intertwiner}
\end{equation}
is an isomorphism of the full quantum
$\mathfrak{gl}_2\oplus\mathfrak{gl}_2$ modules.
To prove this, first apply Lemma
\ref{lem:gct-explicit-string-splitting} to the $V$ generators.
The commuting $W$ generators act on each $V$ highest subspace.
Apply the lemma to that action. The two successive string maps
are exactly \eqref{eq:gct-actual-string-intertwiner}.
The lemma verifies both raising and lowering operators, and
the weights verify all four original diagonal generators and
both total degree-fifteen central characters. This proves the
intertwiner and its bijectivity. The character comparison now
gives $\dim_K H_{ab}=g_{ab}$; no choice of purported highest
source basis vectors was made.

Here is an explicit inverse requiring only the original operators.
For either colour, put
\[
 C=(q-q^{-1})^2FE+qK+q^{-1}K^{-1}.
\]
Its commutators with $K,E,F$ are zero. For example its commutator
with $E$ is
\[
 -(q-q^{-1})E(q^2K-q^{-2}K^{-1})
 +q(q^2-1)EK+q^{-1}(q^{-2}-1)EK^{-1}=0;
\]
the $F$ calculation follows by the same defining relations.
On a highest vector of weight $n$ it acts by
$c_n=q^{n+1}+q^{-n-1}$, hence by that scalar on the whole string.
All source weights $h$ are odd. Consequently the operator
\begin{equation}
 B=FE+\sum_h[(h+1)/2]^2\,\mathrm{pr}_h
       =\frac{C-2}{(q-q^{-1})^2}
 \label{eq:gct-integral-casimir}
\end{equation}
preserves $X_A$. The equality holds because
$q^{h+1}+q^{-h-1}-2=(q^{(h+1)/2}-q^{-(h+1)/2})^2$.
Here $\mathrm{pr}_h$ is projection to the original weight
subspace of the source weight basis, so it is $A$-linear.
The apparent quotient in \eqref{eq:gct-integral-casimir} has
already been replaced by its exact integral operator formula;
it is not a limit used to discard any term.

For $n\in\mathcal N=\{1,3,5,7,9\}$ set
$\beta_n=[(n+1)/2]^2$ and
\begin{equation}
 \Pi_n(B)=\prod_{t\in\mathcal N\setminus\{n\}}
 \frac{B-\beta_t}{\beta_n-\beta_t}.
 \label{eq:gct-source-isotypic-projector}
\end{equation}
The denominators are nonzero since at $q=1$ the five eigenvalues
are respectively $1,4,9,16,25$. The string decomposition proves
that these are orthogonal idempotents with sum one, acting as
identity exactly on the strings of highest ratio weight $n$.

For an original vector $x$ of weight $(u,15-u;v,15-v)$ set
$i=a-u$, $j=b-v$. Keep a pair $(a,b)$ exactly when
$0\leq i\leq n$ and $0\leq j\leq m$. Its inverse component is
\begin{equation}
 e_i\otimes f_j\otimes
 \frac{E_V^{(i)}E_W^{(j)}
       \Pi_n(B_V)\Pi_m(B_W)x}
 {\displaystyle \genfrac{[}{]}{0pt}{}{n}{i}_q
                   \genfrac{[}{]}{0pt}{}{m}{j}_q},
 \qquad \genfrac{[}{]}{0pt}{}{n}{i}_q=\frac{[n]!}{[i]![n-i]!}.
 \label{eq:gct-source-intertwiner-inverse}
\end{equation}
To check both compositions, apply this formula to
$F_V^{(r)}F_W^{(s)}z$ from each string in
\eqref{eq:gct-actual-string-intertwiner}. The projectors kill all
different highest-weight pairs. Weight equality forces $i=r,j=s$
in the surviving pair, and repeated raising gives
$E_V^{(i)}F_V^{(i)}z=\genfrac{[}{]}{0pt}{}{n}{i}_qz$ and its $W$ analogue.
Thus the formula recovers $z$ and every component, proving both
inverse identities. This is a formula on the actual source
generators, with concrete finite products; it is not a character-only
identification.

\subsection{The precise integral localization and specialization}

Let $S$ be the multiplicative set generated by $[1],\ldots,[9]$
and $\beta_n-\beta_t$ for distinct $n,t\in\mathcal N$,
and write $R=S^{-1}A$. The exact Laurent identity
\[
 [r]^2-[s]^2=[r-s][r+s]
\]
follows by expanding the four powers after multiplication by
$(q-q^{-1})^2$, including its sign when $r<s$.
For distinct $r,s\in\{1,2,3,4,5\}$ its two nonzero factors
are signed members of $[1],\ldots,[9]$. Thus this displayed
localization equals $A[([1]\cdots[9])^{-1}]$; the Casimir
denominators were retained and have now been proved to be units.
The lattice $X_A$ remains the original
lattice; this explicitly stated scalar extension is needed for
the following splitting. Every operator in
\eqref{eq:gct-actual-string-intertwiner} and
\eqref{eq:gct-source-intertwiner-inverse} is defined over $R$:
the divided generators already preserve $X_A$, the integral
$B$ does also, and the remaining denominators are units in $R$.
Put $H_{ab,R}=X_R{}_{(a,15-a;b,15-b)}\cap\ker E_V\cap\ker E_W$.
The formulas give mutually inverse $R$-linear maps
\begin{equation}
 \bigoplus_{a,b}L_{(a,15-a),R}\otimes
 L_{(b,15-b),R}\otimes H_{ab,R}\ \cong\ X_R.
 \label{eq:gct-localized-lattice-splitting}
\end{equation}
Indeed the numerator in the inverse formula lies in this highest
kernel over $K$ by the proved inverse, and its annihilation by
both $E$ holds over $R$ because $X_R$ is torsion-free. Both
composition identities also descend from $K$ by torsion-freeness.
Moreover
$H_{ab,R}=\Pi_n(B_V)\Pi_m(B_W)X_R{}_{(a,15-a;b,15-b)}$:
one inclusion follows from the inverse formula at $i=j=0$,
and a highest vector has the indicated two Casimir eigenvalues,
which proves the other inclusion. It is therefore a finite
projective $R$-module of rank $g_{ab}$.

Evaluation at $q=1$ sends $[j]$ to $j$ and the distinct
Casimir differences to differences of $1,4,9,16,25$.
These nonzero differences have prime factors only in
$\{2,3,5,7\}$, and $[1],\ldots,[9]$ already require all four
primes. The universal property of localization therefore gives
\[
 R/(q-1)R=\mathbb Z[1/210].
\]
Tensoring the two inverse maps in
\eqref{eq:gct-localized-lattice-splitting} with this ring keeps
them inverse. In particular the specialized highest kernel is
exactly the specialization of $H_{ab,R}$. One can check this
also directly on the strings: the coefficient of every
non-top raising step is one of $1,\ldots,9$, a unit in
$\mathbb Z[1/210]$, so simultaneous raising kills exactly
the two top positions. No additional highest vectors appear.
The specialized $B$ is
$FE+\sum_h((h+1)/2)^2\mathrm{pr}_h$ and retains its five
distinct eigenvalues. Thus these projectors specialize without
a pole, and the string intertwiners become the classical
divided-power intertwiners. After tensoring with $\mathbb Q$,
the source map $\sigma$ above relates them to the classical
Schur module. This does not claim that the splitting holds
over unlocalized $A$, or identify an integral Schur lattice
without its own lattice comparison.

\subsection{Exact locations of the three retained weights}

The weight $(12,3;9,6)$ has dimension two. Its two components
are $(a,b)=(12,10)$ at positions $(i,j)=(0,1)$ and
$(12,9)$ at positions $(0,0)$, each with multiplicity one.
The complete locations at the next two weights are as follows;
``copies'' is the dimension of the indicated actual $H_{ab}$.
\begin{center}\small
\begin{tabular}{c|cc|c|c}
\toprule
$(u,v)$ & $a$ & $b$ & $(i,j)$ & copies\\\midrule
$(11,9)$&12&10&$(1,1)$&1\\
&12&9&$(1,0)$&1\\
&11&11&$(0,2)$&1\\
&11&10&$(0,1)$&2\\
&11&9&$(0,0)$&3\\\midrule
$(10,9)$&12&10&$(2,1)$&1\\
&12&9&$(2,0)$&1\\
&11&11&$(1,2)$&1\\
&11&10&$(1,1)$&2\\
&11&9&$(1,0)$&3\\
&10&12&$(0,3)$&1\\
&10&11&$(0,2)$&2\\
&10&10&$(0,1)$&4\\
&10&9&$(0,0)$&4\\\bottomrule
\end{tabular}
\end{center}
The corresponding dimensions are $8$ and $19$. Every location
is obtained by $i=a-u$, $j=b-v$ from an actual source isotypic
projector and the proved string isomorphism. A specific original
honest vector such as $T,T_1,T_2,Q$ is located inside these
summands by applying \eqref{eq:gct-source-intertwiner-inverse}
to that vector. Merely knowing its weight does not assign it
to one summand. The original four-vector negative-Borel interval
is still the specified descendant quotient of this full
$1260$-dimensional module; these character and projector formulas
do not identify that interval with an invariant full quantum
group module.

For these four particular vectors the full original source matrices
now permit a stronger determination than their weights alone.
In their original order and with their original scalars they are
\begin{align*}
 T&=(123)(123)(124)(12)(12)(1)(1),\\
 T_1&=(123)(123)(124)(23)(12)(1)(1),\\
 T_2&=-[2]^{-1}(123)(123)(123)(12)(12)(4)(1),\\
 Q&=-[2]^{-1}(123)(123)(124)(13)(12)(4)(1).
\end{align*}
The source's integral height-two commutation relation moves
the columns $(23)(12)$ of $T_1$ to $(12)(23)$ and
the columns $(13)(12)$ of $Q$ to $(12)(13)$ with coefficient
$+1$. Their resulting canonical positive representatives,
including $-[2]^{-1}$ for $T_2,Q$, are basis indices
$126,146,8,144$ of the full original matrix certificate.
The quotient projection on the original ambient words proves
these four equalities, not merely a match of their weights.
The graphical raising rule pairs a $2$ with a later $1$ in
the original $V$ or $W$ word and sums over unpaired $2$'s.
For the four representatives its unpaired-word data are
\[
\begin{array}{c|cc}
 &V&W\\\hline
T&1^9&1^3\\T_1&1^7&1^3\\T_2&1^7&1^3\\Q&1^5&1^3
\end{array}
\]
Thus each raising sum is empty, proving
$E_VT=E_WT=E_VT_1=E_WT_1=E_VT_2=E_WT_2=E_VQ=E_WQ=0$
directly in the original module. Hence
\[
 T\in H_{12,9},\qquad T_1,T_2\in H_{11,9},\qquad
 Q\in H_{10,9}.
\]
The two middle vectors are independent because they are distinct
members of the original honest basis. In particular they span
two dimensions in the three-dimensional actual $H_{11,9}$.
The exact projector computation on all four original vectors
returns the vector itself for this listed pair and zero for
every other pair; its raising and lowering inverse checks use
the full original matrices over $\mathbb Q(q)$.
For example $F_V^{(2)}T$ remains in the isotypic summand
$(\lambda,\mu)=((12,3),(9,6))$, whereas $Q$ is in the different
summand $((10,5),(9,6))$. The original nonzero coefficient of
$Q$ in the honest-basis expansion of $F_V^{(2)}T$ is therefore
not an isotypic projection. Applying $\Pi_5(B_V)\Pi_3(B_W)$
to the complete expansion gives zero: the contribution of the
$Q$ basis term is canceled by the projections of the other
terms. This identity follows because the projectors commute
with $F_V$ and kill $T$ for that highest pair. It preserves
the original coefficient and explains its exact relation to
the full irreducible decomposition.
